\begin{statement}{C.2.UN.67}{Proved}{A torsion-normal elliptic boundary cannot lie on a rational component}
Let \(f:U\to\Delta\) be a proper degeneration of connected
complex two-tori, with smooth total threefold and full
reduced SNC central fibre. Let \(P\) be a rational
component, with no minimality assumption. It cannot
have a smooth elliptic boundary curve \(J\) whose normal
line \(N_{J/P}\) is torsion and which is disjoint from
every other boundary curve on \(P\).

In particular a rational component, including a rational
elliptic surface and any point blowups of a rational
surface, cannot be adjacent to an elliptic Hopf component
or a nonlinear Hopf component. No Hopf leaf can have a
rational parent. The assertion permits rational double
curves and triple points elsewhere in the central fibre.
It imposes neither a zero-Euler hypothesis nor a
minimality hypothesis on the other components.

The obstruction is obtained on the actual component.
Its torsion normal produces the connected elliptic
pencil \(|mJ|\), where \(m=\operatorname{ord}N_{J/P}\).
The actual limiting translation fields force this
pencil to have identity integral monodromy. After its
vertical exceptional curves are contracted, it has
only smooth multiple fibres and coherent Euler
characteristic zero, contradicting
\(\chi(\OO_P)=1\).
\end{statement}

\begin{proof}
\par\medskip\noindent\textbf{The exact pencil made by the torsion normal.}
Suppose that such a component and curve exist, and put
\(N=N_{J/P}\), with exact order \(m\ge1\).
The degree of \(N\) is zero, hence \(J^2=0\).
For every \(k\ge1\) use the actual divisor restriction
sequence
\[
 0\longrightarrow\OO_P((k-1)J)
 \longrightarrow\OO_P(kJ)
 \longrightarrow N^{\otimes k}\longrightarrow0.
 \tag{C.RH.1}
\]
For \(1\le k<m\), the last line is a nontrivial
degree-zero line on the elliptic curve. It has
both \(H^0\) and \(H^1\) zero: a nonzero section
would be nowhere zero, and Serre duality gives
the \(H^1\) assertion. Since \(P\) is rational,
\(H^1(P,\OO_P)=0\) and \(H^0(P,\OO_P)=\C\).
Induction in the actual long exact sequences gives
\[
 h^0(P,\OO_P(kJ))=1,\qquad
 H^1(P,\OO_P(kJ))=0\quad(0\le k<m).
 \tag{C.RH.2}
\]
At \(k=m\), the restriction in (C.RH.1) is
onto \(H^0(J,\OO_J)=\C\). Thus
\[
 h^0(P,\OO_P(mJ))=2,\qquad
 H^0(P,\OO_P(mJ))\longrightarrow H^0(J,\OO_J)
                         \text{ is onto}.
 \tag{C.RH.3}
\]
One section has divisor \(mJ\); a second restricts
to a nonzero constant on \(J\). Their common
zero locus is empty: on \(J\) the second
does not vanish, and off \(J\) the first
does not vanish. They define the actual
holomorphic map
\[
                   q_0:P\longrightarrow\PP^1,\qquad
                             q_0^{-1}(\infty)=mJ .
 \tag{C.RH.4}
\]

We check its connectedness through the Stein map,
rather than assuming that a pencil has connected
fibres. Write \(P\xrightarrow q B\xrightarrow\nu\PP^1\)
for its Stein factor, and let \(d=\deg\nu\).
The injection \(H^1(B,\OO_B)\to H^1(P,\OO_P)\)
gives \(B\simeq\PP^1\).
Since the support of (C.RH.4) is connected,
\(\nu^{-1}(\infty)\) has just one point.
It is totally ramified of order \(d\), and
the actual fibre of \(q\) there is
\((m/d)J\). Its principal local equation
trivializes \(\OO_P((m/d)J)|_J\), so
\[
                   N^{\otimes(m/d)}=\OO_J.
 \tag{C.RH.5}
\]
Exact order \(m\) forces \(d=1\).
Thus \(q=q_0\) has connected fibres.
Adjunction gives \(K_P\cdot J=0\);
a general member \(F\sim mJ\) is smooth
by generic smoothness and satisfies
\[
                 2g(F)-2=F\cdot(K_P+F)=0.
 \tag{C.RH.6}
\]
It is a smooth connected elliptic curve.
This constructs the required elliptic
map on the original nonminimal surface.

Every other boundary component \(C\) on
\(P\) is disjoint from \(J\) by hypothesis.
Consequently \(C\cdot q^*(\infty)=mC\cdot J=0\).
If \(q|_C\) were nonconstant, that intersection
would be its positive degree. Hence every
boundary is vertical for this very map \(q\).

\par\medskip\noindent\textbf{The actual translation fields align this pencil.}
The component has full multiplicity one.
The saturated relative translation module
of UN16 supplies two commuting meromorphic
fields \(V_1,V_2\) on \(P\), independent over
\(\C\), all of whose constant combinations
are semicomplete. Their polar divisors
are contained in the boundary of \(P\).
Indeed they are obtained from holomorphic
ambient relative fields by meromorphic
base coefficients; cancelling the generic
order on \(P\) leaves a holomorphic
restriction at every point of \(P\)
away from the other central components.
This is the original surface and its
actual fields, not a replacement cover.

Remove the finitely many base values
supporting a boundary, a singular fibre
or a vertical zero of the relevant field.
On each remaining compact elliptic fibre
the fields are holomorphic. Their base
coefficients are holomorphic functions
on that fibre and hence are constant.
The connected-fibre map therefore
descends them to commuting meromorphic
fields on \(B=\PP^1\).

These projected fields have no poles.
For clarity, at a pole of order \(a>0\),
the primitive of the one-dimensional
time form has a zero of order \(a+1\).
Continue one inverse branch around
its corresponding time value. Choose
the small loop so that its lifted
base paths avoid every other bad
value. Their compact union has a
compact inverse image under the proper
map \(q\), contained in the holomorphic
locus of the upstairs field. The
upstairs solution therefore continues
along the entire loop. Its final
base value is a different inverse
branch, contradicting the
single-valued maximal solution of
a semicomplete field. This is the
actual continuation map used in
(C.RF.1) and (C.MR.8); it does
not require a minimal surface.

If one projection is nonzero,
the other is a constant multiple,
because the derivative of their
meromorphic ratio along a nonzero
commuting curve field is zero.
Subtract that constant multiple.
The resulting nonzero vertical
field is a translation field on
each compact regular elliptic fibre.
Its commuting transverse field
identifies neighbouring fibres and
preserves its time form. Its two
marked integral periods are
therefore locally constant and,
by analytic continuation on the
connected smooth base, globally
single-valued. Integration of a
nonzero elliptic time form is
injective on \(H_1(F,\Z)\);
the integral monodromy is the
identity.

If both projections vanish,
the fields are vertical translation
fields on each regular fibre.
Their meromorphic ratio is a
nonconstant function pulled
back from \(B\), since the two
fields are independent over
\(\C\). Its connected Stein
fibres are exactly those of
the map \(q\) just constructed.
The actual ambient return germs
(C.RG.4)--(C.RG.12) apply to
this same first integral:
\(\lambda_i=d_i-k_iH\) for its
two continued integral periods.
The right-hand side is
single-valued, and injectivity
of elliptic period integration
again gives identity integral
monodromy. No elliptic pencil
on an ambient neighbourhood
has been assumed.

\par\medskip\noindent\textbf{Relative minimization and the exact singular-fibre comparison.}
Contract only the vertical
\((-1)\)-curves of \(q\).
Each contraction is the
actual smooth surface
blowdown of a rational
exceptional curve, and \(q\)
descends because that curve
is one fibre of the blowdown
and maps to a point of \(B\).
There are finitely many such
steps, since each lowers
\(b_2\). The resulting map
\[
               q_{\min}:P_{\min}\longrightarrow\PP^1
 \tag{C.RH.7}
\]
is relatively minimal. It
has the same marked smooth
fibres and the same identity
monodromy. Moreover
\(\chi(\OO_{P_{\min}})=\chi(\OO_P)=1\),
by the direct-image equality
for every ordinary point
blowup.

We now use the relatively
minimal singular-fibre
classification, precisely in
its stated scope:
\cite[Theorem~6.2 and Table~I]{Kodaira1963II}.
Its nonmultiple non-smooth
types have monodromies
conjugate to nontrivial
unipotent matrices, their
negatives, or the nonidentity
finite-order matrices of
orders \(2,3,4,6\).
Only the smooth type \(I_0\)
has identity monodromy.

For a full multiple fibre
write its local function as
\(t=u h^a\), with \(u\) a
unit and the divisor of
\(h\) primitive. Normalize
the actual base change
\(t=s^a\) by
\[
                         w^a=u,\qquad s=wh .
 \tag{C.RH.8}
\]
The resulting surface map
is \'{e}tale over the
original surface, because
\(w\) is a root of a unit.
Its monodromy is the
\(a\)-th power of the
original one and is
therefore the identity.
It is still relatively
minimal by the following
canonical intersection
comparison. For every
irreducible component \(C\)
of a fibre \(F\) below,
\(F\cdot C=0\) gives
\(C^2\le0\). Adjunction
gives
\(K_{P_{\min}}\cdot C
=2p_a(C)-2-C^2\ge0\),
except possibly for a
smooth rational curve of
square zero or minus one.
The latter is excluded
by relative minimality.
In the square-zero case,
all intersections with
the other fibre components
vanish. Connectedness
would give \(F=bC\),
contradicting
\(K_{P_{\min}}\cdot F=0\).
Thus every component has
nonnegative canonical
intersection. Their
positive weighted sum
is \(K_{P_{\min}}\cdot F=0\),
so each intersection
is zero. For the actual
\'{e}tale covering map
\(\pi\) in (C.RH.8),
\(K_{\rm up}=\pi^*K_{\rm down}\).
Any vertical \((-1)\)-curve
\(E\) upstairs would then
satisfy
\(K_{\rm up}\cdot E
=\deg(E\to C)K_{\rm down}\cdot C=0\),
contrary to adjunction,
which gives minus one.
This uses the cycle
pushforward and does
not assume that the
image curve is embedded
or nonsingular.
Its primitive Kodaira
fibre must consequently
be \(I_0\). Thus every
fibre of \(q_{\min}\)
is a positive multiple
of a smooth elliptic
curve.

The oriented period ratio
is now a single-valued
holomorphic map from a
punctured \(\PP^1\) to
the upper half-plane.
Its bounded Cayley
transform extends across
the missing points and
is constant on \(\PP^1\).
All regular fibres have
one fixed Jacobian.
Their transition maps
are translations, since
the integral monodromy
is the identity. At a
multiple value use
(C.RH.8). The free
central elliptic inertia
is a translation, and
its action on \(H^1(\OO)\)
is trivial. Taking the
actual invariant coefficient
functions gives
\[
 (q_{\min})_*\OO_{P_{\min}}=\OO_{\PP^1},
 \qquad R^1(q_{\min})_*\OO_{P_{\min}}=\OO_{\PP^1}.
 \tag{C.RH.9}
\]
The coherent Leray
sequence now yields
\(\chi(\OO_{P_{\min}})=1-1=0\),
contradicting the preserved
value one in (C.RH.7).
This proves the obstruction
without applying a
minimal-only assertion
to the original \(P\).

\par\medskip\noindent\textbf{The Hopf interface.}
On an elliptic Hopf
component every curve
is a fibre support and
has normal of finite
exact order, by the
actual normalized-root
map in UN37. A nonlinear
Hopf has its one axis,
whose normal has exact
resonance order, including
the secondary comparison
of UN50. Distinct Hopf
curves are disjoint.
There can therefore be
no triple point on an
interface with a Hopf
component: such a point
would make two different
boundary curves meet
inside that Hopf surface.
The actual \(xy=t\)
normal-product map gives
\[
                     N_{J/P}\otimes N_{J/H}=\OO_J.
 \tag{C.RH.10}
\]
It makes \(N_{J/P}\)
torsion, and the same
absence of triple points
makes \(J\) disjoint
from every other
boundary curve on \(P\).
The obstruction just
proved applies.

Every terminal Hopf
pair has one of these
torsion-normal axes,
by the exact
one-boundary canonical
and character comparison
in UN53. The two
nontorsion axes of
an algebraic-dimension-zero
diagonal Hopf cannot
form such a leaf.
This proves all the
asserted parent and
interface exclusions.
\end{proof}


% C.2.RAT.5 from part_C2_rational_triple.tex

